\appendix
\setcounter{equation}{0}
\renewcommand{\theequation}{A.\arabic{equation}}

\chapter*{附录}
\addcontentsline{toc}{chapter}{附录}
\setcounter{chapter}{1}

\def\sectionAname{数学原理}
\section*{附录 A：\sectionAname}
\addcontentsline{toc}{section}{附录 A：\sectionAname}
\setcounter{section}{2}

\subsection*{A.1 属性约简}

\subsubsection*{决策表与不可区分关系}

设决策表为 $DT=(U,C\cup D,V,f)$，其中：\(U\)：对象集合，即样本集；\(C\)：条件属性集合，即特征集；\(D\)：决策属性集合，即标签；\(V\)：属性值域；\(f:U\times(C\cup D)\to V\)：信息函数。

\begin{definition}[不可区分关系]
	设 $B\subseteq C$ 是条件属性子集，则称
	\[
	IND(B)=\{(x,y)\in U\times U\mid \forall a\in B,\ f(x,a)=f(y,a)\}
	\]
	为 $B$ 上的不可区分关系。
\end{definition}

\begin{proposition}
	条件属性子集上的不可区分关系是等价关系。
\end{proposition}

\begin{proof}
	
\end{proof}

\begin{notation}
	由 \(IND(B)\) 诱导的等价类记为：$[x]_B=\{y\in U\mid (x,y)\in IND(B)\}.$ 所有等价类构成划分 $U/IND(B)=\{[x]_B\mid x\in U\}$.
\end{notation}

\subsubsection*{下近似、正域与依赖度}

\begin{definition}[下近似]
	设 $U$ 的子集 $X$, 则称 
	\[
	\underline{B}(X)=\{x\in U\mid [x]_B\subseteq X\}
	\]
	为 $X$ 关于 $B$ 的一个下近似。
\end{definition}

\begin{notation}
	决策属性 \(D\) 将论域划分为决策类 $U/D=\{D_1,D_2,\dots,D_r\}$
\end{notation}

\begin{definition}[正域]
	条件属性子集 \(B\) 关于决策属性 \(D\) 的正域为：
	\[
	POS_B(D)=\bigcup_{i=1}^{r}\underline{B}(D_i)
	\]
\end{definition}

\begin{definition}[依赖度]
	定义条件属性 $B$ 关于决策属性 $D$ 的依赖度为：
	\[
	\gamma_B(D)=\frac{|POS_B(D)|}{|U|},
	\]
	表示用条件属性 \(B\) 能够确定分类的对象比例。
\end{definition}

\subsubsection*{约简与核}

\begin{definition}[约简]
	称 \(R\) 为 \(C\) 相对于 \(D\) 的一个约简，若 \(R\subseteq C\) 满足 $\gamma_R(D)=\gamma_C(D)$ 且对任意 \(a\in R\)，都有  
	\[
	\gamma_{R\setminus\{a\}}(D)<\gamma_R(D).
	\]
\end{definition}

\begin{definition}[核]
	称所有约简的交集为核，记为
	\[
	Core(C)=\bigcap\{R\mid R\text{ 是约简}\}
	\]
\end{definition}

\begin{proposition}
	任取 $c\in Core(C),$ 都有 $\gamma_{C\setminus\{c\}}(D) < \gamma_C(D)$.
\end{proposition}
\begin{proof}
	
\end{proof}

\subsubsection*{属性重要度}
\begin{definition}
	定义属性 \(a\) 相对于当前属性子集 \(R\) 的重要度为：
	\[
	Sig(a,R,D)=\gamma_{R\cup\{a\}}(D)-\gamma_R(D).
	\]
	重要度越大，说明加入属性 \(a\) 后，正域增加越多，分类能力提升越明显。
\end{definition}